\chapter{Convolution}

In the previous chapter we made a single, far-reaching discovery: if a system
is \emph{linear} and \emph{time-invariant} (an LTI system), then its entire
behavior is locked down by one curve --- the way it responds to a single sharp
poke, the impulse response. Everything the system will ever do to any input,
however complicated, is already encoded in that one response.

This chapter turns that principle into a concrete calculation. We are going to derive the exact
machine that turns ``input plus impulse response'' into ``output.'' That
machine is \textbf{convolution}. Filtering is convolution. Echo and reverb are
convolution. Blurring and sharpening images is convolution. Smoothing a noisy
accelerometer trace is convolution. The Fourier transform's central operational
theorem is a statement about convolution. And when we eventually smooth,
debounce, and extract features from the signals coming off a wrist sensor, every one
of those steps is a convolution.

Our promise for this chapter: you will not memorize the convolution formula.
You will \emph{see} it. We build it from the geometry of decomposing a signal
into impulses, watch it run as a flip-shift-multiply-sum animation rendered in
stills, prove its properties from that picture, and run it on real examples ---
a moving-average smoother and a digital echo --- in both the continuous and
discrete worlds.

\chapterline

\section{The one idea: build any signal out of impulses}
\label{ch13:decompose}

Everything flows from a clean geometric observation. Take a discrete-time
signal $x[n]$ --- just a sequence of numbers, one per sample index. Here is a
representation that might initially look like mere bookkeeping, but actually
unlocks the entire operational theory:

\begin{equation}
\label{ch13:sifting-discrete}
x[n] \;=\; \sum_{k=-\infty}^{\infty} x[k]\,\delta[n-k].
\end{equation}

Read this slowly. The unit impulse $\delta[n]$ is $1$ at $n=0$ and $0$
everywhere else. So $\delta[n-k]$ is a single spike sitting at position $k$.
Multiply that spike by the number $x[k]$, and you get a spike of height $x[k]$
at position $k$. Now add up one such scaled spike for every $k$. What do you
get back? Exactly the original signal. At index $n$, only the $k=n$ term in the
sum survives (every other $\delta[n-k]$ is zero there), and it contributes
$x[n]$. The equation is a tautology --- but a \emph{generative} one.

\begin{intuitionbox}
A signal is a stack of impulses. Picture $x[n]$ as a bar chart. Each bar is one
impulse: it lives at a particular time, and it has a particular height. The
whole signal is nothing more than that collection of weighted, shifted spikes,
laid side by side. We have not changed the signal at all --- we have only
chosen to \emph{see} it as a sum of elementary pieces. Choosing the right
elementary pieces is the recurring trick of this entire book. With sinusoids we
got the Fourier transform. With impulses we are about to get convolution.
\end{intuitionbox}

\begin{figure}[H]
\centering
\begin{tikzpicture}
\begin{axis}[
  width=12cm, height=5cm,
  xlabel={$n$}, ylabel={$x[n]$},
  xmin=-0.7, xmax=6.7, ymin=0, ymax=3.4,
  xtick={0,1,2,3,4,5,6}, ytick={0,1,2,3},
  axis lines=left, ymajorgrids, grid style={gray!18},
]
\addplot+[ycomb, mark=*, thick, color=sigblue, mark options={fill=sigblue}]
  coordinates {(0,0.5)(1,1.5)(2,3)(3,2)(4,1)(5,0.5)(6,0.2)};
\end{axis}
\end{tikzpicture}
\caption{A discrete signal is a row of impulses. Each stem is one term
$x[k]\,\delta[n-k]$ of Eq.~\eqref{ch13:sifting-discrete}: a spike at position
$k$ with height $x[k]$. The signal \emph{is} the collection.}
\label{ch13:fig-stems}
\end{figure}

Why bother rewriting a signal as a sum of impulses? Because of what an LTI
system does with a sum of impulses. This is where linearity and
time-invariance work together.

\section{From impulse response to convolution}
\label{ch13:derive}

Let $T\{\cdot\}$ denote our LTI system: feed it an input, it produces an
output. Define its \textbf{impulse response} as the output when the input is a
single unit impulse at the origin:
\begin{equation}
\label{ch13:imp-resp-def}
h[n] \;\coloneqq\; T\{\delta[n]\}.
\end{equation}
This is the one curve from the previous chapter. Now watch the system handle an
arbitrary input $x[n]$, using the two properties in turn.

\textbf{Step 1 --- time-invariance.} If a poke at time $0$ produces $h[n]$,
then a poke at time $k$ produces the same shape, just shifted to start at $k$:
\[
T\{\delta[n-k]\} \;=\; h[n-k].
\]
The system does not care \emph{when} you poke it; it only cares about the delay
between poke and observation.

\textbf{Step 2 --- linearity.} The input is a sum of scaled, shifted impulses,
Eq.~\eqref{ch13:sifting-discrete}. A linear system lets us push $T$ through the
sum and through the scalars: the response to a sum is the sum of responses, and
the response to a scaled input is the scaled response. So
\begin{align}
y[n] \;=\; T\{x[n]\}
   &= T\!\left\{ \sum_{k} x[k]\,\delta[n-k] \right\} \notag\\
   &= \sum_{k} x[k]\, T\{\delta[n-k]\}
   \qquad\text{(linearity)} \notag\\
   &= \sum_{k} x[k]\, h[n-k].
   \qquad\text{(time-invariance)}
\label{ch13:conv-discrete-derived}
\end{align}

That last line \emph{is} the convolution sum. We did not invent it; we were
forced into it. Once you accept that a signal is a stack of impulses, and that
the system responds to each impulse with a shifted copy of $h$, the only thing
left to do is add up all those copies, each weighted by how strong its impulse
was. That sum is convolution.

\begin{happeningbox}
Convolution is a superposition of echoes. Every input sample $x[k]$ throws a
scaled copy of the impulse response $h$ into the future (and present), landing
at time $k$. At any output time $n$, you stand still and collect the
contributions from \emph{all} those copies that happen to reach $n$. A copy
launched at time $k$ contributes $x[k]\,h[n-k]$ --- its launch strength $x[k]$
times the impulse response evaluated at the elapsed time $n-k$. The output is
the running total of overlapping echoes. Hold onto this ``superposition of
shifted, scaled copies'' picture; it is the truest mental model of
convolution.
\end{happeningbox}

\begin{definition}[Discrete convolution]
\label{ch13:def-discrete}
The convolution of two discrete-time signals $x$ and $h$ is the signal
\begin{equation}
(x \conv h)[n] \;=\; \sum_{k=-\infty}^{\infty} x[k]\,h[n-k].
\end{equation}
We also write $y = x \conv h$. The output of an LTI system is the convolution
of its input with its impulse response: $y = x \conv h$.
\end{definition}

The continuous-time story is identical, with the sum becoming an integral. A
continuous signal is a ``sum'' of infinitely many infinitely thin impulses, the
sifting property of the Dirac delta:
\begin{equation}
\label{ch13:sifting-cont}
x(t) \;=\; \int_{-\infty}^{\infty} x(\tau)\,\delta(t-\tau)\,\dd\tau .
\end{equation}
Run the same two-step argument --- time-invariance turns $\delta(t-\tau)$ into
$h(t-\tau)$, linearity pushes $T$ inside the integral --- and you land on the
continuous convolution.

\begin{definition}[Continuous convolution]
\label{ch13:def-cont}
\begin{equation}
(x \conv h)(t) \;=\; \int_{-\infty}^{\infty} x(\tau)\,h(t-\tau)\,\dd\tau .
\end{equation}
\end{definition}

\begin{connectionbox}
Equation~\eqref{ch13:sifting-discrete} is the impulse-basis cousin of writing a
vector in coordinates. In the linear-algebra companion, a vector
$\vec{x} = \sum_k x_k\,\vec{e}_k$ is a weighted sum of standard basis vectors
$\vec{e}_k$. Here the ``basis vectors'' are shifted impulses $\delta[n-k]$, and
$x[k]$ are the coordinates. An LTI system is a linear map; once you know what it
does to each basis element ($\delta[n-k] \mapsto h[n-k]$), linearity fixes what
it does to everything. Convolution is just that linear map written out in the
impulse basis. The companion's inner-product view returns later: each output
sample $y[n]$ is the inner product of the input with a flipped, shifted copy of
$h$.
\end{connectionbox}

\section{Flip, shift, multiply, sum --- the moving picture}
\label{ch13:flipshift}

Definition~\ref{ch13:def-discrete} has an $h[n-k]$ in it, and that minus sign in
front of $k$ governs the geometry of convolution. Let us decode it as a
recipe you can run by hand. To compute the single output value $y[n]$ at one
chosen index $n$:

\begin{enumerate}
\item \textbf{Flip.} Take $h[k]$ and reflect it about $k=0$ to get $h[-k]$. The
impulse response is now running backwards.
\item \textbf{Shift.} Slide the flipped response right by $n$ to get $h[n-k]$.
It is now anchored at the output time $n$.
\item \textbf{Multiply.} For every $k$, multiply $x[k]$ by the overlapping value
$h[n-k]$.
\item \textbf{Sum.} Add up all those products. The total is $y[n]$.
\end{enumerate}

Then move to the next output index, shift the flipped response one more step,
and repeat. The flipped response \emph{slides across} the input; at each
position you take the running dot product of the overlap. That sliding-overlap
motion is convolution made visible.

\begin{whybox}
Why flip? The minus sign is not arbitrary, and it is not the same as
correlation (which uses $h[n+k]$ and does \emph{not} flip). It comes straight
from the echo picture. An impulse launched at time $k$ has been traveling for
$n-k$ samples by the time we observe at $n$. Early launches ($k$ small) have
been traveling longest, so they show us the \emph{late} part of $h$; recent
launches show us the early part of $h$. When you line the responses up against
the input axis $k$, that ``oldest launch sees latest $h$'' bookkeeping reverses
the order of $h$ --- it flips. The flip is causality and elapsed-time
bookkeeping made geometric.
\end{whybox}

Let us run the machine on a tiny example we can check by hand, then draw two
frames of the slide.

\begin{example}[Convolving two short sequences]
\label{ch13:ex-byhand}
Let $x = [\,1,\ 2,\ 3\,]$ at indices $n=0,1,2$, and let
$h = [\,1,\ 1\,]$ at indices $n=0,1$ (a two-tap summer). The output has length
$3+2-1 = 4$. Using $y[n] = \sum_k x[k]\,h[n-k]$:
\begin{align*}
y[0] &= x[0]h[0] = 1\cdot 1 = 1, \\
y[1] &= x[0]h[1] + x[1]h[0] = 1\cdot 1 + 2\cdot 1 = 3, \\
y[2] &= x[1]h[1] + x[2]h[0] = 2\cdot 1 + 3\cdot 1 = 5, \\
y[3] &= x[2]h[1] = 3\cdot 1 = 3.
\end{align*}
So $x \conv h = [\,1,\ 3,\ 5,\ 3\,]$. Notice the output is longer than either
input: convolution \emph{spreads} a signal out. Notice too that the same
numbers appear if you treat it as polynomial multiplication: $(1 + 2z + 3z^2)
(1 + z) = 1 + 3z + 5z^2 + 3z^3$. That is not a coincidence; we return to it in
Section~\ref{ch13:polynomial}.
\end{example}

\begin{figure}[H]
\centering
\begin{tikzpicture}
\begin{axis}[
title style={font=\small},
  name=f0,
  width=12cm, height=3.3cm,
  title={Frame 0: computing $y[0]$ --- leading edge contact ($n=0$)},
  title style={font=\small},
  xlabel={$k$}, xmin=-2.7, xmax=4.7, ymin=0, ymax=3.6,
  xtick={-2,-1,0,1,2,3,4}, ytick={0,1,2,3},
  axis lines=left, ymajorgrids, grid style={gray!15},
]
\addplot+[ycomb, mark=*, thick, color=sigblue, mark options={fill=sigblue}]
  coordinates {(0,1)(1,2)(2,3)};
\addplot+[ycomb, mark=square*, thick, color=specred, mark options={fill=specred}]
  coordinates {(-1,1)(0,1)};
\node[above, font=\footnotesize, text=targetgreen] at (axis cs:0,1.2) {overlap: $1 \times 1 = 1$};
\end{axis}

\begin{axis}[
title style={font=\small},
  name=f1, at={(f0.below south west)}, anchor=above north west, yshift=-0.45cm,
  width=12cm, height=3.3cm,
  title={Frame 1: computing $y[1]$ --- flipped $h$ slid right to $n=1$},
  title style={font=\small},
  xlabel={$k$}, xmin=-2.7, xmax=4.7, ymin=0, ymax=3.6,
  xtick={-2,-1,0,1,2,3,4}, ytick={0,1,2,3},
  axis lines=left, ymajorgrids, grid style={gray!15},
]
\addplot+[ycomb, mark=*, thick, color=sigblue, mark options={fill=sigblue}]
  coordinates {(0,1)(1,2)(2,3)};
\addplot+[ycomb, mark=square*, thick, color=specred, mark options={fill=specred}]
  coordinates {(0,1)(1,1)};
\node[above, font=\footnotesize, text=targetgreen] at (axis cs:0.5,2.1) {overlap: $1\times 1 + 2\times 1 = 3$};
\end{axis}

\begin{axis}[
title style={font=\small},
  name=f2, at={(f1.below south west)}, anchor=above north west, yshift=-0.45cm,
  width=12cm, height=3.3cm,
  title={Frame 2: computing $y[2]$ --- flipped $h$ slid right to $n=2$ (peak)},
  title style={font=\small},
  xlabel={$k$}, xmin=-2.7, xmax=4.7, ymin=0, ymax=3.6,
  xtick={-2,-1,0,1,2,3,4}, ytick={0,1,2,3},
  axis lines=left, ymajorgrids, grid style={gray!15},
]
\addplot+[ycomb, mark=*, thick, color=sigblue, mark options={fill=sigblue}]
  coordinates {(0,1)(1,2)(2,3)};
\addplot+[ycomb, mark=square*, thick, color=specred, mark options={fill=specred}]
  coordinates {(1,1)(2,1)};
\node[above, font=\footnotesize, text=targetgreen] at (axis cs:1.5,3.1) {overlap: $2\times 1 + 3\times 1 = 5$};
\end{axis}
\end{tikzpicture}
\caption{Three stills of the flip-and-drag convolution sequence for Example~\ref{ch13:ex-byhand}. Blue stems
are the input $x[k]$; red squares are the flipped-and-shifted response
$h[n-k]$. To find $y[n]$, multiply the overlapping blue and red values and sum.
Frame 0 ($n=0$): leading edge overlap yields $y[0]=1$. Frame 1 ($n=1$): overlap at $k=0,1$ gives $y[1]=1{+}2=3$.
Frame 2 ($n=2$): overlap at $k=1,2$ gives $y[2]=2{+}3=5$. Marching the red pattern one step right at each clock tick sweeps out the entire response.}
\label{ch13:fig-slide}
\end{figure}

\begin{walkbox}
Before reading on, compute $y[0]$ and $y[3]$ for
Example~\ref{ch13:ex-byhand} by sliding the red pattern in
Figure~\ref{ch13:fig-slide} all the way left (so it overlaps only $x[0]$) and
all the way right (only $x[2]$). You should get $1$ and $3$. When does the slide
start producing nonzero output, and when does it stop? The answer ---
overlap begins when the leading edge of $h$ first touches $x$ and ends when its
trailing edge leaves --- is exactly why the output length is
$\text{len}(x) + \text{len}(h) - 1$.
\end{walkbox}

\section{Worked example: the moving average}
\label{ch13:movavg}

The most common smoother on earth is the moving average, and it is a
convolution. An $M$-point moving average replaces each sample with the average
of itself and its $M-1$ predecessors:
\begin{equation}
\label{ch13:movavg-def}
y[n] \;=\; \frac{1}{M}\sum_{m=0}^{M-1} x[n-m]
       \;=\; (x \conv h)[n], \qquad
h[n] = \begin{cases} 1/M, & 0 \le n \le M-1,\\ 0, & \text{otherwise.}\end{cases}
\end{equation}
The impulse response is a \emph{rectangle} of height $1/M$ and width $M$. Flip a
rectangle and it looks the same; sliding it across the input and summing the
overlap is literally ``average everything under the box.'' That is why the
moving average is sometimes called a boxcar filter.

\begin{figure}[H]
\centering
\begin{tikzpicture}
\begin{axis}[
  width=13cm, height=5.2cm,
  xlabel={$n$}, ylabel={amplitude},
  xmin=0, xmax=60, ymin=-1.8, ymax=1.8,
  axis lines=left, ymajorgrids, grid style={gray!15},
  legend pos=north east, legend style={font=\scriptsize, fill=white},
]
\addplot[color=sigblue!50, thick] coordinates {
(0,0.10)(2,0.55)(4,0.05)(6,0.75)(8,0.35)(10,0.95)(12,0.55)(14,1.05)(16,0.70)
(18,1.10)(20,0.62)(22,0.95)(24,0.40)(26,0.70)(28,0.10)(30,0.30)(32,-0.25)
(34,-0.05)(36,-0.55)(38,-0.30)(40,-0.85)(42,-0.55)(44,-1.05)(46,-0.70)
(48,-1.10)(50,-0.62)(52,-0.95)(54,-0.40)(56,-0.70)(58,-0.20)(60,-0.45)};
\addlegendentry{noisy $x[n]$}
\addplot[color=specred, very thick, smooth] coordinates {
(0,0.35)(4,0.55)(8,0.72)(12,0.86)(16,0.95)(20,0.97)(24,0.90)(28,0.72)
(32,0.45)(36,0.10)(40,-0.28)(44,-0.60)(48,-0.82)(52,-0.92)(56,-0.85)(60,-0.62)};
\addlegendentry{$y[n]=(x\conv h)$, $M{=}7$}
\end{axis}
\end{tikzpicture}
\caption{A moving average is convolution with a rectangle. The jagged blue
trace is a noisy signal; the smooth red trace is the output of convolving it
with a width-$M$ boxcar. Averaging neighbors cancels uncorrelated wiggles while
leaving the slow trend intact. (Traces drawn schematically to show the effect.)}
\label{ch13:fig-movavg}
\end{figure}

\begin{trapbox}
The moving average is a low-pass filter, but a mediocre one. Its frequency
response is the sinc-shaped magnitude $\abs{H(\omega)} = \abs{\sinc(M\omega/2) /
\sinc(\omega/2)}$ (we derive this kind of thing in the next chapter). The trap:
that response has big \emph{sidelobes} and even passes through zero at certain
frequencies while letting others leak. So a boxcar smooths, but it also smears
and ripples in the frequency domain. When you need real attenuation, reach for a
designed filter, not a raw moving average. The moving average wins on one axis
only: it is the cheapest possible low-pass, and with a running sum it costs one
add and one subtract per sample regardless of $M$.
\end{trapbox}

\begin{connectionbox}
Why does averaging kill noise? Lean on the probability companion. If each sample
is signal plus independent zero-mean noise of variance $\sigma^2$, averaging $M$
of them leaves the signal (roughly) intact but shrinks the noise variance to
$\sigma^2/M$ --- the variance of a mean of independent variables. So the
noise standard deviation drops by $\sqrt{M}$. Convolution with a length-$M$
boxcar buys you a $\sqrt{M}$ improvement in signal-to-noise, at the cost of
blurring features narrower than $M$ samples. That tradeoff --- noise rejection
versus time resolution --- is the eternal tension of smoothing.
\end{connectionbox}

\section{Worked example: a digital echo}
\label{ch13:echo}

Echo is the purest demonstration of the superposition-of-copies picture. A
single delayed-and-attenuated repeat is the system
\begin{equation}
\label{ch13:echo-def}
y[n] \;=\; x[n] + a\,x[n-D],
\end{equation}
where $D$ is the delay in samples and $0 < a < 1$ is how much quieter the echo
is. What impulse response produces this? Feed in $\delta[n]$: the output is
$\delta[n] + a\,\delta[n-D]$, so
\[
h[n] = \delta[n] + a\,\delta[n-D]
     = \{\,1 \text{ at } n{=}0,\ \ a \text{ at } n{=}D\,\}.
\]
Two spikes: the direct sound and one echo. Convolving any music $x[n]$ with this
two-spike $h$ lays down a copy of the whole song at full volume, plus a second
copy delayed by $D$ samples and scaled by $a$. That is exactly what an echo
\emph{is} --- a delayed, fainter copy summed with the original. The echo picture
of convolution from Section~\ref{ch13:derive} was not a metaphor; for this $h$
it is literal.

\begin{figure}[H]
\centering
\begin{tikzpicture}
\begin{axis}[
  width=12cm, height=4.4cm,
  xlabel={$n$ (samples)}, ylabel={$h[n]$},
  xmin=-3, xmax=34, ymin=0, ymax=1.25,
  xtick={0,10,20,30}, ytick={0,0.5,1},
  axis lines=left, ymajorgrids, grid style={gray!15},
]
\addplot+[ycomb, mark=*, thick, color=sigblue, mark options={fill=sigblue}]
  coordinates {(0,1)(10,0.6)(20,0.36)(30,0.216)};
\node[anchor=south, font=\scriptsize] at (axis cs:0,1.06) {direct};
\node[anchor=south, font=\scriptsize] at (axis cs:10,0.66) {echo 1};
\node[anchor=south, font=\scriptsize] at (axis cs:20,0.42) {echo 2};
\end{axis}
\end{tikzpicture}
\caption{Impulse response of a feedback echo (here $D=10$, $a=0.6$, with the
output fed back so each echo spawns the next). Each spike is one delayed, fainter
copy. Convolving a sound with this comb of spikes stamps that pattern of repeats
onto every moment of the sound. A single-tap echo would show just the first two
spikes; a reverb's impulse response is a dense, decaying cloud of thousands.}
\label{ch13:fig-echo}
\end{figure}

\begin{applicationbox}
Real reverb is convolution with a measured room. To capture how a cathedral or a
recording studio sounds, engineers play a known test signal (a starter pistol, a
balloon pop, or a swept sine) and record the result; deconvolving the input out
leaves the room's impulse response --- a dense, exponentially decaying spray of
reflections lasting seconds. A \emph{convolution reverb} plugin then convolves
your dry recording with that captured impulse response, transplanting the
cathedral onto your voice. The same idea, run in reverse, lets acousticians
diagnose flutter echoes and resonances in a hall. Convolution lets you place any
dry recording inside the acoustics of a real room.
\end{applicationbox}

\section{Properties and their geometric meaning}
\label{ch13:props}

Convolution obeys a handful of algebraic laws. We could grind them out by
manipulating sums, but the point of this book is to \emph{see} them. Each falls
out of either the echo picture or the slide picture.

\textbf{Commutativity: $x \conv h = h \conv x$.} Substitute $m = n-k$ in the
sum; the roles of input and impulse response swap. Geometrically: it does not
matter whether you flip-and-slide $h$ across $x$ or $x$ across $h$ --- the set
of overlapping products is the same. Practically, this is liberating: ``filter
the signal'' and ``stamp the filter's pattern with the signal'' compute the same
thing.

\textbf{Associativity: $(x \conv h_1) \conv h_2 = x \conv (h_1 \conv h_2)$.}
Cascading two LTI systems is itself an LTI system, and its impulse response is
the convolution of the two individual responses (Figure~\ref{ch13:fig-cascade}).
So a chain of filters can be collapsed into one equivalent filter --- precompute
$h_1 \conv h_2$ once and convolve with it forever after.

\begin{figure}[H]
\centering
\begin{tikzpicture}[
  box/.style={draw, thick, minimum width=1.5cm, minimum height=0.85cm, rounded corners=2pt},
  >=Latex, node distance=0.85cm,
]
% cascade
\node (xi) {$x$};
\node[box, right=of xi] (h1) {$h_1$};
\node[box, right=of h1] (h2) {$h_2$};
\node[right=of h2] (yi) {$y$};
\draw[->] (xi) -- (h1);
\draw[->] (h1) -- (h2);
\draw[->] (h2) -- (yi);
\node[font=\scriptsize, below=0.25cm of h1.south east] {two systems in series};
% equals
\node[right=0.7cm of yi] (eq) {$\equiv$};
% combined
\node[right=0.7cm of eq] (xi2) {$x$};
\node[box, right=of xi2, minimum width=2.0cm] (hc) {$h_1\conv h_2$};
\node[right=of hc] (yi2) {$y$};
\draw[->] (xi2) -- (hc);
\draw[->] (hc) -- (yi2);
\node[font=\scriptsize, below=0.25cm of hc.south] {one equivalent system};
\end{tikzpicture}
\caption{Associativity as a wiring fact. Cascading an $h_1$ filter into an $h_2$
filter is identical to a single filter whose impulse response is $h_1 \conv h_2$.
You can precompute the combined response once and convolve with it forever after.}
\label{ch13:fig-cascade}
\end{figure}

\textbf{Distributivity over addition:
$x \conv (h_1 + h_2) = x \conv h_1 + x \conv h_2$.} Two LTI systems wired in
parallel (same input, outputs summed) equal one system whose impulse response is
the sum of theirs. This is just linearity again.

\textbf{Identity: $x \conv \delta = x$.} Convolving with a single unit spike at
the origin returns the signal unchanged. Of course --- $\delta$ is the
``do-nothing'' impulse response, the wire that passes its input straight
through. And $x \conv \delta[n-D] = x[n-D]$: convolving with a shifted spike is
a pure delay.

\textbf{Shift: if $y = x \conv h$, then $x[n-D] \conv h = y[n-D]$.} Delay the
input, delay the output by the same amount. This is just time-invariance,
restated.

\begin{formalbox}
Commutativity, made precise. Starting from $(x\conv h)[n] = \sum_k x[k]\,h[n-k]$,
let $m = n-k$, so $k = n-m$ and the sum over all $k$ becomes a sum over all $m$:
\[
(x\conv h)[n] = \sum_{m} x[n-m]\,h[m] = \sum_m h[m]\,x[n-m] = (h\conv x)[n].
\]
The same substitution works for the continuous integral. Every other property
admits an equally short proof; try associativity as an exercise by writing the
double sum and reordering.
\end{formalbox}

\begin{aestheticbox}
Notice the algebraic structure here. Under convolution, signals form a
commutative algebra: there is an addition, a commutative and associative
``multiplication'' ($\conv$), it distributes over addition, and there is a
multiplicative identity ($\delta$). Structurally, it obeys the same rules as ordinary
number multiplication --- and as Section~\ref{ch13:polynomial} shows, discrete convolution
has the exact algebra of polynomial multiplication, while the Fourier transform is the
change of coordinates that turns $\conv$ into pointwise multiplication.
\end{aestheticbox}

\section{Convolution is polynomial (and integer) multiplication}
\label{ch13:polynomial}

Example~\ref{ch13:ex-byhand} dropped a hint: convolving the coefficient lists
$[1,2,3]$ and $[1,1]$ gave $[1,3,5,3]$, the coefficients of
$(1+2z+3z^2)(1+z)$. This is a theorem, not a fluke. Encode a sequence $x[n]$ as
the coefficients of a polynomial (or a power series), $X(z) = \sum_n x[n]\,z^n$.
Then multiplying two such polynomials,
\[
X(z)\,H(z) = \Big(\sum_k x[k]z^k\Big)\Big(\sum_m h[m]z^m\Big)
           = \sum_n \Big(\underbrace{\sum_k x[k]\,h[n-k]}_{(x\conv h)[n]}\Big) z^n,
\]
collects, into the coefficient of $z^n$, exactly the convolution sum. The
$z^k \cdot z^m = z^{n}$ bookkeeping ($k+m=n$) is the same as the $n-k$ index in
convolution.

\begin{whybox}
This is why the $z$-transform turns convolution into plain multiplication:
$\Ztr\{x \conv h\} = X(z)\,H(z)$. It is also why the Fourier transform does the
same with $z = e^{\jj\omega}$ --- the convolution theorem says
$\FT\{x\conv h\} = X\cdot H$. Convolution in the time domain, multiplication in
the transform domain: one operation seen in two coordinate systems. Transforming
to the frequency domain turns convolution's flip-and-slide sum into a direct
pointwise product.
\end{whybox}

And the same logic explains why grade-school long multiplication of integers
\emph{is} convolution-with-carries: multiplying $123 \times 11$ produces the
digit-convolution $[1,3,5,3]$ before you carry. Long multiplication without carries
is already discrete convolution.

\begin{historybox}
The word ``convolution'' (from Latin \emph{convolvere}, to roll together) is
younger than the operation itself. Integrals of the form $\int f(\tau)g(t-\tau)\,\dd\tau$
appeared throughout the 18th and 19th centuries in the work of d'Alembert, Euler,
Laplace, and Poisson, solving differential equations and calculating the probability
distributions of summed independent random variables. The German term \emph{Faltung}
(``folding'') was standard well into the 20th century, beautifully capturing the
physical geometric action of flipping one signal backwards before sliding it past another.

When linear system theory crystallized in the mid-20th century, convolution emerged
as the universal law governing every filter, acoustic echo, and transmission channel.
Yet direct convolution carries a steep computational penalty: convolving a signal of
length $N$ with an impulse response of length $M$ requires roughly $NM$ multiplications
and additions. For long audio tracks or real-time radar processing, direct time-domain
convolution was computationally crippling on early hardware.

This bottleneck drove the development of transform methods. As we will explore in
Chapter~14 (\emph{Impulse and Frequency Response}), representing signals as sums of
complex exponentials turns convolution into simple pointwise multiplication. Later,
in Chapter~17 (\emph{Practical Filtering}), we will see how the Fast Fourier Transform
harnesses this duality through overlap-add block processing, slashing the computational
cost from $O(NM)$ down to $O(N \log M)$.
\end{historybox}

\section{Edges and the cost of convolution}
\label{ch13:practical}

Two practical facts bite the moment you convolve real, finite signals.

\textbf{Length and edges.} Convolving a length-$N$ signal with a length-$M$
impulse response gives a length-$(N+M-1)$ result. The extra $M-1$ samples are
the ``tail'' where the impulse response is sliding off the end. Near both edges
the response only partially overlaps the data, so those output samples are based
on fewer terms --- they are less trustworthy. In practice you choose a boundary
policy: zero-pad (assume the signal is $0$ outside, the default and the source of
those tails), repeat the edge value, reflect the signal, or wrap it around
(circular convolution, which the DFT performs). For a causal real-time filter
you keep only the first $N$ outputs and accept a startup transient.

\textbf{Cost.} Direct convolution costs about $N \cdot M$ multiply-adds: for
each of $N$ outputs you sum $M$ products. When both signals are long this is
brutal --- $O(N^2)$. The convolution theorem rescues us: transform both signals
with the FFT ($O(N\log N)$), multiply pointwise, inverse-transform. For long
filters this ``fast convolution'' is dramatically cheaper, and it is how serious
reverb and large-kernel image filters are actually computed.

\begin{lstlisting}[language=Python, caption={Discrete convolution three ways: by hand, with NumPy, and the FFT route.}]
import numpy as np

def conv_direct(x, h):
    N, M = len(x), len(h)
    y = np.zeros(N + M - 1)
    for n in range(N + M - 1):
        for k in range(N):              # flip-shift-multiply-sum
            m = n - k
            if 0 <= m < M:
                y[n] += x[k] * h[m]
    return y

x = np.array([1., 2., 3.])
h = np.array([1., 1.])
print(conv_direct(x, h))                # [1. 3. 5. 3.]
print(np.convolve(x, h))                # [1. 3. 5. 3.]  -- same

# Fast convolution via the FFT (zero-pad to avoid circular wrap):
L = len(x) + len(h) - 1
Y = np.fft.ifft(np.fft.fft(x, L) * np.fft.fft(h, L)).real
print(np.round(Y, 6))                   # [1. 3. 5. 3.]  -- same again
\end{lstlisting}

\begin{buildbox}
On an Arduino-class microcontroller you will convolve with \emph{short} impulse
responses --- a 5-to-15-tap smoother on each IMU axis, say --- and you will do
it sample-by-sample in real time, never storing the whole signal. Keep the last
$M$ samples in a small \emph{circular buffer} and compute one output per new
sample:
\begin{lstlisting}[language=C, caption={Streaming FIR convolution (length-M) in C.}]
#define M 7
float h[M] = {1./M,1./M,1./M,1./M,1./M,1./M,1./M};  // boxcar smoother
float buf[M] = {0};      // circular buffer of recent inputs
int   pos = 0;

float fir_step(float x_new) {
    buf[pos] = x_new;                 // overwrite oldest sample
    float y = 0.0f;
    int idx = pos;
    for (int m = 0; m < M; m++) {     // dot product with the taps
        y += h[m] * buf[idx];
        idx = (idx == 0) ? M-1 : idx-1;
    }
    pos = (pos + 1) % M;
    return y;
}
\end{lstlisting}
This is $O(M)$ per sample and uses $O(M)$ memory --- trivial for $M$ in the
single digits at IMU sample rates of a few hundred hertz. For a boxcar
specifically you can do even better with a running sum (one add, one subtract),
but the FIR form above generalizes to any designed filter.
\end{buildbox}

\begin{goalbox}
Convolution is the first real signal-processing step in the wrist tracker. Raw
accelerometer and gyroscope streams are noisy and jittery; before you can detect
a gesture you smooth each axis with a short low-pass convolution (a boxcar or,
better, a designed FIR) to kill sensor hash and hand tremor. Tap and flick
detection then convolves the smoothed signal with a short \emph{matched} kernel
shaped like the gesture's acceleration bump --- a sliding template match, which
(Section~\ref{ch13:flipshift}) is convolution with a flipped template. Even
debouncing a contact event is a tiny convolution that spreads and thresholds the
pulse.

Be clear about what this buys you. Smoothing and template-convolution on a
wrist IMU detect \emph{coarse, discrete} motions --- a flick, a tap, a wrist
roll, a double-tap. They do \emph{not} recover continuous individual finger
joint angles; the wrist simply does not move enough, distinctly enough, when
only a finger curls. Per-finger sensing needs EMG or force myography on the
forearm, and continuous finger \emph{pose} needs sensors on the fingers
themselves (rings or a glove). Those signals, too, get convolved --- EMG is
band-pass filtered and its envelope is extracted by convolution --- but the
sensor, not the math, sets the ceiling on what gestures you can read.
Convolution cleans and matches the signal you have; it cannot conjure
information the sensor never captured.
\end{goalbox}

\section{A glance ahead: 2D convolution for images}
\label{ch13:twod}

Everything here generalizes to two dimensions with a second sum (or integral).
For an image $x[m,n]$ and a small kernel $h[i,j]$,
\begin{equation}
\label{ch13:conv2d}
(x \conv h)[m,n] \;=\; \sum_{i}\sum_{j} x[i,j]\,h[m-i,\,n-j].
\end{equation}
The recipe is the same: flip the kernel (now in both axes), slide it across the
image, and at each position take the sum of the overlapping products. A
$3\times3$ averaging kernel blurs; a kernel with a positive center and negative
surround sharpens or finds edges. We devote Chapters~30--32 to it, because image
filtering, edge detection, and the convolutional layers of neural networks are
all this one operation in 2D. The intuition you built in 1D --- a sliding,
weighted, flipped overlap --- transfers wholesale; only the bookkeeping gains an
index.

\begin{recoverybox}
Putting the chapter together: a signal is a stack of impulses; an LTI system
answers each impulse with a shifted, scaled copy of its impulse response; the
output is the superposition of all those copies; written as a formula that
superposition is the convolution sum $\sum_k x[k]h[n-k]$; computed by hand it is
flip, shift, multiply, sum; it is commutative, associative, distributive, with
$\delta$ as identity; it has the exact algebra of polynomial multiplication, which is why a
transform turns it into a pointwise product; and it is the foundation behind
filtering, echo, smoothing, and --- in 2D --- image processing.
\end{recoverybox}

\begin{connectionbox}
\textbf{Capabilities, Downsides, and Forward Roadmap}
\begin{itemize}
  \item \textbf{What this unlocked}: Convolution provides the exact mathematical engine for all linear
  time-invariant systems. Whether smoothing noisy sensor traces, simulating concert-hall acoustics,
  or blurring image pixels, convolution translates the system's impulse response into the complete output.
  \item \textbf{Honest downsides and limitations}: Time-domain convolution hides the spectral character
  of the filter. Looking at the numbers of $h[n]$ alone, it is difficult to immediately tell which frequencies
  pass, which are attenuated, or where resonance occurs. In addition, calculating long convolutions in the
  time domain costs $\mathcal{O}(NM)$ operations.
  \item \textbf{Where we go next}: Chapter~14 introduces \textbf{Impulse and Frequency Response}. We show that
  complex exponentials are the eigenfunctions of convolution, turning time-domain sums into
  frequency-domain multiplications $Y(\omega) = X(\omega)H(\omega)$ that directly expose resonances,
  bandwidth, and filter design.
\end{itemize}
\end{connectionbox}

\section{Exercises}
\label{ch13:exercises}

\begin{exercise}
\label{ch13:ex1}
By hand, compute $x \conv h$ for $x = [\,2,\ 0,\ -1,\ 3\,]$ (indices $0$--$3$)
and $h = [\,1,\ 2\,]$ (indices $0$--$1$). State the length of the result and
verify it by multiplying the polynomials $2 - z^2 + 3z^3$ and $1 + 2z$.
\end{exercise}

\begin{exercise}
\label{ch13:ex2}
A three-point moving average has $h = [\,1/3,\ 1/3,\ 1/3\,]$. Convolve it with
the unit step $x[n] = 1$ for $n \ge 0$ (and $0$ before). Sketch the output for
$n = -1$ through $n = 4$ and describe the startup transient. How many samples
until the output reaches its steady value?
\end{exercise}

\begin{exercise}
\label{ch13:ex3}
Prove associativity of convolution, $(x \conv h_1) \conv h_2 =
x \conv (h_1 \conv h_2)$, by writing the left side as a double sum and
reindexing. Then interpret the result as a statement about cascaded LTI systems.
\end{exercise}

\begin{exercise}
\label{ch13:ex4}
A single-echo system has $h[n] = \delta[n] + 0.5\,\delta[n-4]$. (a) Write the
difference equation relating $y[n]$ to $x[n]$. (b) Convolve $h$ with itself.
What system does $h \conv h$ represent, and how many echoes does it have? (c) In
general, what does convolving an echo with itself $k$ times do?
\end{exercise}

\begin{exercise}
\label{ch13:ex5}
Show that $x \conv \delta[n-D] = x[n-D]$ directly from the convolution sum.
Explain in one sentence why this means ``convolving with a shifted spike is a
pure delay.''
\end{exercise}

\begin{exercise}
\label{ch13:ex6}
The flip in convolution is what distinguishes it from cross-correlation,
$r[n] = \sum_k x[k]\,h[k+n]$ (no flip). (a) For $h$ symmetric, $h[n] = h[-n]$,
show convolution and correlation give the same answer. (b) Give a short,
non-symmetric $h$ where they differ, and compute both for some input. (c) Which
one is the right tool for sliding-template gesture matching, and why does the
flip matter there?
\end{exercise}

\begin{exercise}
\label{ch13:ex7}
You convolve a length-$1000$ signal with a length-$64$ filter. (a) How many
multiply-adds does direct convolution take, roughly? (b) Estimate the cost of
the FFT route ($\sim 3 \cdot L\log_2 L$ for three transforms of length $L$, plus
$L$ multiplies, with $L = 1024$). (c) At what filter length does the FFT route
start to win, qualitatively?
\end{exercise}

\begin{exercise}
\label{ch13:ex8}
Using the connection to the probability companion, suppose each sample of a flat
signal is corrupted by independent noise of standard deviation $\sigma = 0.2$.
You smooth with an $M$-point moving average. (a) What is the noise standard
deviation of the output? (b) How large must $M$ be to cut the noise to
$0.05$? (c) What feature width (in samples) do you lose by choosing that $M$,
and how is this the noise-versus-resolution tradeoff?
\end{exercise}

\begin{exercise}
\label{ch13:ex9}
Write (in pseudocode or Python) a streaming convolution that processes one
sample at a time with a length-$M$ FIR filter using a circular buffer, matching
the C sketch in this chapter. Test it against \texttt{numpy.convolve} on a
random input and confirm the outputs agree (allowing for the edge/startup
samples).
\end{exercise}

\begin{exercise}
\label{ch13:ex10}
Take the $2\times2$ image $X = \left[\begin{smallmatrix} 1 & 2 \\ 3 & 4
\end{smallmatrix}\right]$ and the $2\times2$ kernel $H =
\left[\begin{smallmatrix} 0 & 1 \\ 1 & 0 \end{smallmatrix}\right]$. Compute the
full 2D convolution $X \conv H$ using Eq.~\eqref{ch13:conv2d} (remember to flip
$H$ in both axes). What size is the result, and which output entry corresponds to
the kernel fully overlapping the image?
\end{exercise}
